\documentclass[10pt,a4paper]{article}
\usepackage[margin=1.7cm]{geometry}
\usepackage{enumitem}\usepackage{titlesec}\usepackage[hidelinks]{hyperref}\usepackage{parskip}
\titleformat{\section}{\normalfont\large\bfseries}{\thesection}{0.6em}{}
\titlespacing*{\section}{0pt}{6pt}{2pt}
\setlist[itemize]{topsep=2pt,itemsep=1pt,parsep=0pt,leftmargin=1.3em}
\pagestyle{empty}
\begin{document}
\begin{center}
{\Large\bfseries IMU-Based Camera Stabilisation: Tilt and Pan Control}\\[2pt]
{\large Lab 2 \textendash{} open-loop control with a BNO055 IMU and two servos}\\[4pt]
{\small Embedded Real-Time Control laboratory \textendash{} M.Sc. Control Systems Engineering, University of Padova}\\
{\small Author: Kiavash Rabbani Soudabi (ID 2141822) \textendash{} Course held by Prof.\ Alberto Morato \textendash{} 2025}
\end{center}

\section*{Purpose}
Implement open-loop control on the T-Bot: use an inertial measurement unit to stabilise the tilt of a camera mount
and to drive its pan axis, then improve the behaviour with two bonus tasks. The IMU is a BNO055 read over
I\textsuperscript{2}C; the two axes are actuated with PWM servos.

\section*{What I implemented}
\begin{itemize}
  \item \textbf{Exercise 1 -- tilt stabilisation from the accelerometer:} the tilt angle is estimated from the
        measured Y acceleration, $\theta=\arcsin(a_y/g)$ converted to degrees, and the camera command is the negative
        of that angle, clamped to $\pm45^\circ$.
  \item \textbf{Bonus 1 -- smooth pan control from the gyroscope:} the pan command is a pure proportional law on the
        Z-axis angular rate, $u_{pan}=-k_p\,\omega_z$ with $k_p=30$ (the value I tuned during the lab), also clamped
        to $\pm45^\circ$.
  \item \textbf{Actuator mapping:} each angle is mapped to the servo duty as $150+\theta\cdot(50/45)$ counts and
        passed through a \texttt{saturate()} function that enforces the servo limits, so the mechanical range is
        never exceeded.
  \item \textbf{Telemetry:} I restructured the datalogger record to carry the interesting signals
        (\texttt{tilt\_angle}, \texttt{y\_accelerate}, \texttt{pan\_control}, \texttt{z\_gyro}) and fixed the order of
        the logger update/send calls so that the transmitted packets are complete; data was acquired with the MATLAB
        serial/WiFi datalogger for analysis.
  \item \textbf{Bonus 2 -- sensor fusion:} combining accelerometer and gyroscope information for a more stable
        tilt/pan response.
\end{itemize}

\section*{Results}
The measured step responses for the exercise and for both bonuses are included as figures in the report: the tilt
loop holds the camera level for moderate platform inclination, and the gyro-based pan follows the rotation rate far
more smoothly than the accelerometer-only solution. The report documents the control laws, the clamping strategy and
the measurement results.

\section*{Repository contents}
\texttt{report/Embedded-Report-lab2.pdf} (submitted report, LaTeX source and figures included) $\cdot$
\texttt{src/user\_code\_main.c} (USER CODE sections of my \texttt{main.c}; course framework not redistributed) $\cdot$
\texttt{docs/} (this overview).

\section*{Skills demonstrated}
IMU integration over I\textsuperscript{2}C (BNO055), inclinometer estimation, proportional control design and gain
tuning, accelerometer/gyroscope sensor fusion, PWM servo actuation and saturation handling, fixed-rate real-time
loops, measurement and analysis with MATLAB.

\end{document}
